% Pursuing Stacks (cotes 134-2 to 134-8): archive cote 134-5, site page 49 (PDF page 50), typescript page 236.
% typescript. Draft compared: ps4.tex:2418-2499.
% AI-transcribed by Claude Opus 5.5 (claude-opus-5-5, adaptive thinking, effort high), Message Batches API, batch ps2 (msgbatch_01NVwr1Yuo8qJf5rykeGBXPK), 2026-11-12; reviewed against both supplied scan halves; see CHECK_DECISIONS.tsv.
\scanpage{134-5/49}
\mspage{236}

To sum up some of the main relationships between the three asphericity
notions just introduced ($M_{as}$ , $W$, $M_{as}$-aspheric functors), \over{L}{l}et's state
one more tautological proposition, which is very much a paraphrase of
the display given earlier (prop. on p. 214) of the manifold aspects of
the notion of an aspheric map between small categories:

\hand{\uline{\textcolor{black}{Proposition}}} 4. Let $(M,M_{as})$ be an asphericity structure, $A$ a small
category, $i\colon A\hand{\to} M$ a functor, factoring through $M_{as}$ . Consider the
following conditions on $i$ :

(i) $i$ is $M_{as}$-aspheric (or \over[typed]{\illegible{1}}{a} ``testing functor'' for $M_{as}$), i.e. \xout{forx}
\[ M_{as} = (i^{*})^{-1}(A^{\wedge}_{as}) \qquad . \]
(i') For $x$ in $M_{as}$ , $i^{*}(x)$ is aspheric, i.e.
\[ M_{as} \hand{\subset} (i^{*})^{-1}(A^{\wedge}_{as}) \quad . \]
(i'') (Here $B$ is a given small full subcategory \inserted[typed]{of $M_{as}$} \inserted{containing $i(A)$ and} which``\hand{\uline{\textcolor{black}{generates}}}'' the
asphericity structure $M_{as}$ , namely such that the inclusion functor
$j\colon B\hand{\hookrightarrow} M$ is $M_{as}$-aspheric i.e. $M_{as} = (j^{*})^{-1}(B^{\wedge}_{as})$.) For any $x$ in $B$,
$i^{*}(x)$ is aspheric, i.e.
\[ B \hand{\subset} (i^{*})^{-1}(A^{\wedge}_{as}) \quad . \]
(ii) For any map $u$ in $M$, $u$ is \xout{aspheric} \inserted[typed]{a weak equivalence} iff $i^{*}(u)$ is, i.e.
\[ W = (i^{*})^{-1}(\underline{W}_{A}) \quad . \]
(ii') If the map $u$ in $M$ is \xout{aspheric,} \inserted[typed]{a weak equivalence} so is $i^{*}(u)$ , i.e.
\[ W \hand{\subset} (i^{*})^{-1}(\underline{W}_{A}) \quad . \]
(ii'') (Here $B$ is given as in (i'') above) For any map \unsure{u} in $B$, $i^{*}(u)$
is a weak equivalence, \inserted{i.e. $Fl(B) \hand{\subset} (i^{*})^{-1}(\underline{W}_{A})$}. \margin{——}

% A small upright mark above The remains unidentified; the margin strokes are retained.
\margin{——} The conditions (i)(i')(i'') are equivalent and imply all others, and
we have the tautological implications (ii)$\hand{\Rightarrow}$(ii')$\hand{\Rightarrow}$(ii''). If $M$ admits
a final object $e_M$ and \inserted[typed]{if A and} the asphericity structure \inserted{$M_{as}$} \deleted{is} \inserted[typed]{are} aspheric, then \xout{all} \inserted{all}
\inserted[typed]{six} \xout{\illegible{1}} conditions \inserted[typed]{except the last} are equivalent; and all \inserted{six} are equivalent if \inserted[typed]{moreover} \xout{x} $e_M \hand{\in} \mathrm{Ob}\, B$ .

Proof. The imp\xout{p}lications (i)$\hand{\Rightarrow}$(i')$\hand{\Rightarrow}$(i'') are tautological, on
the other hand (i'') just means that the induced functor $f\colon A\hand{\to} B$ is as-%
pheric, which by prop.1 a) implies that $i$ is aspheric. On the other
hand (i)$\hand{\Rightarrow}$(ii) by the definition of $W$ (cf. prop. 2). If $e_M$ exists
and \inserted{A and} the asphericity structure $M_{as}$ \over{is}{are} aspheric, and if $B$ contains $e_M$ ,
then (ii'') implies that the maps $x\hand{\to} e_M$ for $x$ in $B$ \xout{xxxxx} are trans-%
formed by $i^{*}$ into a weak equivalence, and as $e_{A^{\wedge}}$ is aspheric this
implies $i^{*}(x)$ is aspheric i.e. (i''), which \xout{impl} proves the last state-%

\margin{49}
